\section{A conditional cylinder sieve}\label{sec:cylinder-sieve}

The restrictions $n+b_j\in Y$ still involve the residue coordinates
outside the fixed set $\mathcal F$.  We replace their product by a sum
of conditions from short lists of primitive specifications.  In each
term the remaining coordinates are averaged before taking an absolute
value; this preserves the cancellation needed in the trace estimate.
The error will contain many specifications with distinct private prime
coordinates, which supply enough arithmetic constraints to apply
\Cref{lem:triangular-saving}.

The intersection rank and the extraction of constraints from its
witnesses adapt Pilatte's construction
\cite[Definition~13.3 and Lemmas~13.4 and~13.6]{Pilatte26}.
The associated truncation belongs to the composite-modulus sieve
developed in \cite[Section~3]{HR21} and
\cite[Appendix~B]{Pilatte26}.  We prove the finite-cylinder statement
and its application here, including the effect of conditioning.

\subsection{Truncating an intersection poset}

Let $\Omega=\prod_{p\in P}\Omega_p$ be a finite product, with
$|\Omega_p|\ge2$ for every $p$.  An \emph{exact cylinder} specifies one
value at each coordinate in a subset of $P$ and imposes no condition at
the other coordinates.  Its \emph{support} is the set of specified
coordinates, and its \emph{width} is the size of that support.

Let $\mathscr E$ be a finite family of proper, nonempty exact cylinders.
Identify duplicate cylinders.  Let $\mathscr P$ consist of the nonempty
intersections of subfamilies of $\mathscr E$, including the empty-family
intersection $\widehat0=\Omega$.  Order $\mathscr P$ by reverse inclusion:
$I\le J$ means $I\supseteq J$.  For $I\in\mathscr P$, define
\begin{equation}\label{eq:cylinder-rank}
 \operatorname{rk}(I)=
 \max\left\{|\mathscr A|:\begin{array}{l}
   \mathscr A\subseteq\{E\in\mathscr E:I\subseteq E\},\\
   \text{each }E\in\mathscr A\text{ has a support coordinate}\\
   \text{absent from every other member of }\mathscr A
 \end{array}\right\}.
\end{equation}
The empty family is allowed in this maximum.  A coordinate as in
\eqref{eq:cylinder-rank} is called \emph{private} to its member of
$\mathscr A$.  All events in such a family have compatible prescribed
values because they contain the nonempty intersection $I$.

Write $\mu(\widehat0,I)$ for the M\"obius coefficients of this finite
poset, characterized by
\[
 \sum_{I\le J}\mu(\widehat0,I)=\1_{J=\widehat0}.
\]
Their role is inclusion--exclusion with equal intersections collected
into one term.  This is M\"obius inversion in the incidence algebra of a
finite poset; see \cite[Section~3]{Rota1964}.  The chain expansion used
below is the one in \cite[Section~3, Proposition~6]{Rota1964}.

\begin{lemma}[Finite-cylinder truncation]\label{lem:cylinder-truncation}
Suppose that the cylinders in $\mathscr E$ have width at most $w\ge1$.
For an integer $t\ge1$, put $M=tw$, and let
\[
 \mathscr H_t=\{J\in\mathscr P:\operatorname{rk}(J)\ge t,
       \ J\text{ is generated by at most }t\text{ members of }\mathscr E\}.
\]
Then the rank is nondecreasing on $\mathscr P$, every $I\in\mathscr P$
is generated by at most $\operatorname{rk}(I)$ events, and
\begin{equation}\label{eq:cylinder-width}
 |\operatorname{supp}(I)|\le w\operatorname{rk}(I).
\end{equation}
For every $x\in\Omega$,
\begin{equation}\label{eq:cylinder-truncation}
 \left|\1_{x\notin\bigcup_{E\in\mathscr E}E}
       -\sum_{\operatorname{rk}(I)<t}\mu(\widehat0,I)\1_I(x)\right|
 \le 2^M\max(1,M^{M+1})\sum_{J\in\mathscr H_t}\1_J(x).
\end{equation}
Every coefficient retained in the sum on the left has absolute value
at most $\max(1,M^{M+1})$.  Every intersection on the right has width
at most $M$, even if its rank is greater than $t$.
\end{lemma}

\begin{proof}
If $I\le J$, every event containing $I$ also contains $J$.  Thus every
family admitted for the maximum defining $\operatorname{rk}(I)$ is
admitted for $\operatorname{rk}(J)$, proving monotonicity.
Choose an inclusion-minimal family generating $I$.  Each member must
have a private coordinate: otherwise all its prescribed values already
follow from the other members, since the prescriptions are compatible,
and it can be removed.  Its size is therefore at most
$\operatorname{rk}(I)$, proving \eqref{eq:cylinder-width} as well.

Fix $J\in\mathscr P$ and let $m=|\operatorname{supp}(J)|$.  An element
of the interval $[\widehat0,J]$ is determined by its support: its values
must be the corresponding restrictions of those prescribed by $J$.
Consequently this interval has at most $2^m$ elements.  Along a strict
chain from $\widehat0$ to $J$, supports grow strictly, so there are at
most $m$ steps.  A chain of $k$ steps is specified by recording, for
each coordinate of $\operatorname{supp}(J)$, the step at which it
enters.  There are at most $k^m$ such chains.  Expanding the inverse of
the poset's upper triangular incidence matrix, or applying the M\"obius
recursion repeatedly, gives the alternating sum over these chains.
Thus, for $m\ge1$,
\begin{equation}\label{eq:cylinder-mobius}
 |\mu(\widehat0,J)|\le\sum_{k=1}^m k^m\le m^{m+1};
\end{equation}
the coefficient at $\widehat0$ is $1$.  This proves the asserted
coefficient bound for intersections of rank below $t$.

For $x\in\Omega$, let $J_x$ be the intersection of all events satisfied
at $x$, with $J_x=\widehat0$ if none are satisfied.  An intersection
$I\in\mathscr P$ contains $x$ if and only if $I\le J_x$.  Indeed any
family generating $I$ then consists of events satisfied at $x$.
It follows that
\begin{equation}\label{eq:cylinder-full-sum}
 \sum_{I\in\mathscr P}\mu(\widehat0,I)\1_I(x)
    =\sum_{I\le J_x}\mu(\widehat0,I)
    =\1_{J_x=\widehat0}.
\end{equation}
Since no forbidden event is the whole space, the last expression is
precisely the avoidance indicator in \eqref{eq:cylinder-truncation}.

If $\operatorname{rk}(J_x)<t$, monotonicity shows that the truncated
sum contains the entire interval $[\widehat0,J_x]$, so it is exact.
Suppose instead that $\operatorname{rk}(J_x)\ge t$.  Given any
$I\le J_x$ of rank below $t$, successively intersect it with events
satisfied at $x$ until reaching $J_x$.  Let $K$ be the first intersection
whose rank is at least $t$, and let $K^-$ be its immediate predecessor.
Replace the generating family for $K^-$ by an inclusion-minimal one.
It has at most $t-1$ members.  Adjoining the event that produces $K$
shows that $K$ is generated by at most $t$ events.  Thus
\[
       I\le K\le J_x,\qquad K\in\mathscr H_t,
       \qquad |\operatorname{supp}(K)|\le M.
\]
This argument allows the rank to jump by more than one.  The
re-minimization of the predecessor, rather than the length of the
successive list, bounds the number of generators.

Every satisfied low-rank $I$ is therefore below a satisfied member of
$\mathscr H_t$.  For any one such member there are at most $2^M$
possible predecessors $I$, by the interval bound already proved.
The exact avoidance indicator is zero in the present case.  Bounding
each retained coefficient by \eqref{eq:cylinder-mobius} now gives
\eqref{eq:cylinder-truncation}.
\end{proof}

\subsection{Applying the truncation after conditioning}

Return to a fixed numerical line $\boldsymbol d$.  Its fixed prime set
$\mathcal F$ consists of its core labels and its repeated labels, as in
\Cref{sec:thinning}.  In particular, core primes absent from the line
are not added to $\mathcal F$.  Write
$R_{\boldsymbol d}(n_{\mathcal F})$ for the indicator of the conditions
retained in \Cref{prop:retained-lines}.

Condition on $n_{\mathcal F}$ with $R_{\boldsymbol d}=1$.  At every
$n+b_j$, take the primitive forbidden specifications compatible with
these fixed residues.  Omit their fixed-coordinate tests, leaving
exact cylinders on the coordinates $\mathcal S\setminus\mathcal F$.
Each has width at most $LJ+1$.  There is no whole-space event: such an
event would be a qualifying primitive specification using only fixed
primes, excluded by the retained conditions.  Avoiding these residual
events is exactly the condition $n+b_j\in Y$ for every $j$.

For a list $\mathfrak L$ of primitive specifications attached at line
indices $0,\ldots,\ell$, let $I_{\mathfrak L}(n)$ be the indicator that
all its specifications qualify, including their fixed-coordinate
tests.  The empty list has indicator $1$.  Let
$\mathscr L_{<t_0}(\boldsymbol d)$ denote all ordered such lists of length
less than $t_0$.  These are finite families at the fixed graph
parameters.  We will use the bound
\begin{equation}\label{eq:cylinder-budget}
 M_0=t_0(LJ+1)=O(B^{1-\rho+4\eps}\log B),\qquad
 M_0\log B=o(B).
\end{equation}
In particular, the coefficient and error factors in
\Cref{lem:cylinder-truncation} are $\exp(O(M_0\log B))$.

A low-rank intersection is generated by fewer than $t_0$ residual
events.  Choose one realizing primitive specification for each event.
The conjunction of their full qualification indicators equals the
intersection indicator at the conditioned residues.  Different
intersections cannot have the same chosen generating list, since that
list determines its intersection.  We may therefore bound their sum
by the sum over $\mathscr L_{<t_0}(\boldsymbol d)$ without a multiplicity
loss.

For an intersection in $\mathscr H_{t_0}$, choose at most $t_0$
generating events and exactly $t_0$ events witnessing its rank.  The
latter may be selected from a larger irredundant family; deleting other
members preserves their private coordinates.  The witnessing events
contain the intersection.  Consequently adjoining them to the
generating list leaves its indicator unchanged.  After choosing
realizing primitive specifications, each witness has a prime outside
$\mathcal F$ that occurs in none of the other witness specifications.
These choices may depend on $n_{\mathcal F}$, which is harmless:
the preceding assertions and the uniform coefficient bounds apply
pointwise at each fixed-coordinate assignment.

We have thus reduced the truncation error to jointly qualifying lists
of at most $2t_0$ specifications, including $t_0$ witnesses with private
primes outside $\mathcal F$.  The next lemma shows that these error
lists impose many triangular constraints on their numerical prime
labels.

\begin{lemma}[Constraints from private primes]\label{lem:cylinder-witnesses}
Let $\boldsymbol d$ be a numerical closed line, and let $\mathcal F$ contain
its core prime labels and all labels occurring at least twice on the
line.  Suppose $t_0$ primitive specifications attached at line vertices
qualify simultaneously.  Suppose that each has a prime outside
$\mathcal F$ absent from every other specification in this family.
For sufficiently large $B$, their numerical labels satisfy a triangular
system on at least $B^{2\eps}$ distinct selected prime variables of the
kind in \Cref{lem:triangular-saving}.  Each constraint uses one intrinsic
suffix displacement, or the difference of two witness prefixes and a
line segment.  Its nonzero expressions have logarithmic size
$O(B\log B)$ even when numerical label-size restrictions are dropped.
\end{lemma}

\begin{proof}
Order the witnesses by nondecreasing attachment indices
$k_1,\ldots,k_{t_0}$; break ties arbitrarily.  For witness $j$, choose a
private prime $q_j\notin\mathcal F$, write $p_j$ for its extra prime,
and choose a tail prime $r_j$ as in
\Cref{lem:primitive-properties}.  Thus $r_j$ occurs only in its final
constant block.  That lemma supplies two coefficient facts:
\begin{enumerate}[label=(\roman*)]
\item the intrinsic suffix displacement divisible by $p_j$ is linear
in a tail-only prime, with coefficient nonzero modulo $p_j$;
\item if a prime occurs before the constant tail, its coefficient in
the prefix displacement immediately preceding that tail is nonzero
modulo $r_j$.
\end{enumerate}
All private primes $q_j$ are distinct.  Each occurs at most once as a
label on the original line, because it is outside $\mathcal F$.
We divide into three cases to control every dependence on another
selected variable.

\paragraph{Private primes in intrinsic suffix constraints.}
Suppose at least $t_0/10$ witnesses have $q_j=p_j$ or have $q_j$
occurring only in the constant tail.  In the first situation $q_j$
divides the intrinsic suffix displacement, which is nonzero by
simplicity and does not involve $q_j$.  In the second situation use
that suffix constraint as a linear congruence in $q_j$ modulo $p_j$;
its coefficient is nonzero by (i).  The private-prime property excludes
every other selected $q$ from the witness, including from its modulus.
These constraints have no dependencies on the other selected
variables and hence form a triangular system.

\paragraph{Tail primes first or last appearing among the witnesses.}
Suppose at least $t_0/10$ witnesses have $r_j$ absent from every earlier
witness.  Select these tail primes in increasing witness order.  They
are distinct: equality between two would make the latter prime occur
in an earlier witness.  Use the intrinsic suffix congruence for each
selected $r_j$.  No later selected tail prime occurs in this witness,
by its defining absence property.  This also excludes later selected
variables from the modulus $p_j$.  Thus each constraint uses only its
own variable, earlier selected variables, and nonselected variables,
and its coefficient is nonzero modulo $p_j$ by (i).
If instead at least $t_0/10$ witnesses have $r_j$ absent from every later
witness, the same argument in decreasing witness order applies.

\paragraph{Comparing two activities of a shared tail prime.}
If none of these cases applies, fewer than $3t_0/10$ witnesses have
been excluded by their three respective conditions.  Retain at least
$t_0/2$ witnesses for which $q_j$ occurs before the tail and $r_j$
appears in both an earlier and a later witness.  Since $q_j$ occurs
at most once on the original line, it is absent either from all steps
$1,\ldots,k_j$ or from all steps $k_j+1,\ldots,\ell$.  At least
$t_0/4$ of the retained witnesses satisfy the same one of these two
absence conditions.

First consider the family absent from steps $1,\ldots,k_j$, ordered
increasingly by attachment index.  For its witness $j$, choose an
earlier witness $i=i(j)$ using $r_j$.  In witness $i$, choose a vertex
at which $r_j$ is active: its starting vertex if $r_j$ is its extra
prime, or the origin of an edge carrying $r_j$ otherwise.  Let $T_i$
be this vertex's offset relative to that witness's start.  Let $P_j$
be the prefix displacement of witness $j$ immediately preceding its
constant tail.  Joint qualification implies
\begin{equation}\label{eq:cylinder-comparison}
 b_{k_j}-b_{k_i}+P_j-T_i\equiv0\pmod{r_j}.
\end{equation}
The coefficient of $q_j$ in this congruence is exactly its coefficient
in $P_j$: privateness excludes it from witness $i$, and the selected
absence condition excludes it from the line segment between
$k_i$ and $k_j$.  This coefficient is nonzero modulo $r_j$ by (ii).

Now take a later selected private variable $q_s$.  It occurs in
neither witness $i$ nor witness $j$, by privateness.  Moreover
$k_s\ge k_j$, and its own absence condition excludes it from every
line step through $k_s$, hence from the segment in
\eqref{eq:cylinder-comparison}.  Thus no later selected variable occurs
in that congruence.  The modulus $r_j$ is not any selected private
prime: it differs from $q_j$ because $q_j$ occurs before the tail,
and it differs from all other $q_s$ by privateness.  This proves
triangularity in increasing order.  Equal attachment indices merely
make the corresponding line segment empty.

For the family absent from steps after $k_j$, order the witnesses
decreasingly and compare witness $j$ to a later witness containing
$r_j$.  Equation \eqref{eq:cylinder-comparison} has the same form, now
with its line segment after $k_j$.  A later selected variable in this
decreasing order has attachment $k_s\le k_j$ and is absent from all
steps after $k_s$, so it is absent from that segment.  Privateness
excludes it from the two witness prefixes.  Fact (ii) again supplies
the nonzero coefficient.  This gives a triangular system in decreasing
order.  The attained comparison displacement may equal zero; that
causes no difficulty, since we use an invertible linear congruence,
not a divisibility constraint with the selected prime as modulus.

Every case provides at least $t_0/20$ selected variables for sufficiently
large $B$, allowing for integer parts.  This is greater than
$B^{2\eps}$.  Each label is a squarefree product of at most $J$ primes
at most $e^B$, so its logarithm is at most $BJ=O(B\log B)$ even without
the divisor-bin restrictions.  A suffix or a comparison above uses
$O(\ell+L)$ such terms.  Every nonzero expression consequently has
logarithmic size $O(B\log B)$.  The constraint choices are specified
by witness and line positions and prime slots, within the enumeration
allowed in \Cref{lem:triangular-saving}.  All its hypotheses are now
verified.
\end{proof}

\begin{proposition}[Reduction to short qualification lists]
\label{prop:cylinder-reduction}
For the retained line and fixed-coordinate data of
\Cref{prop:retained-lines}, there is a factor
$\mathfrak C_B=\exp(O(M_0\log B))$ such that
\begin{align}
 &\sum_{\boldsymbol d}\E_{\mathcal F}
   \left[R_{\boldsymbol d}
     \left|\E_{\mathcal S\setminus\mathcal F}
          K_{\boldsymbol d}(n)\prod_{j=0}^{\ell}\1_Y(n+b_j)\right|\right]
       \notag\\
 &\quad\le \mathfrak C_B\sum_{\boldsymbol d}\E_{\mathcal F}
   \left[R_{\boldsymbol d}
     \sum_{\mathfrak L\in\mathscr L_{<t_0}(\boldsymbol d)}
       \left|\E_{\mathcal S\setminus\mathcal F}
                      K_{\boldsymbol d}(n)I_{\mathfrak L}(n)\right|\right]
       +O\!\left(e^{-B^{1+\eps/2}}\right).
       \label{eq:cylinder-reduction}
\end{align}
Here each sum runs over numerical closed lines, $\mathcal F$ is the
fixed prime set of that line, and every list indicator imposes full
qualification of its specifications.
\end{proposition}

\begin{proof}
At fixed $\boldsymbol d,n_{\mathcal F}$ with $R_{\boldsymbol d}=1$, apply
\Cref{lem:cylinder-truncation} to the residual cylinders with
$t=t_0$ and $w=LJ+1$.  Multiply the identity with its pointwise error
bound by $K_{\boldsymbol d}$ and integrate the residual coordinates.
For the low-rank sum, take the absolute value after each such integral
and use the generating lists already constructed.  The coefficient
bound gives the main term on the right of
\eqref{eq:cylinder-reduction}.

For the error, take the absolute value of $K_{\boldsymbol d}$ inside the
integral and use the generating and witnessing lists described above.
After integrating the fixed coordinates, their indicators require
joint qualification of at most $2t_0$ specifications, including the
$t_0$ witnesses with private primes outside $\mathcal F$.
We may drop the retained-data restriction in this positive bound.
\Cref{lem:cylinder-witnesses} supplies at least $B^{2\eps}$ triangular
constraints for each nonzero contribution.

Refine the error count by the lit or unlit status of every line
occurrence.  Joint qualification fixes one residue for each distinct
prime used by the specifications, or is inconsistent and contributes
zero.  For each refinement, \Cref{lem:crude-count} gives
$A^{\ell J}\prod_p p^{-1}$, with the product over distinct prime labels,
rather than charging a shared residue repeatedly.  The line and lists
have at most
\[
       \ell J+2t_0(LJ+1)=O(B\log B)
\]
prime slots.  Their equality patterns, occurrence statuses and all
choices of witness constraints have cost $\exp(O(B\log^2 B))$.
These are exactly the
positive majorant and description budget required by
\Cref{lem:triangular-saving}.  Its reverse elimination applies to the
system just supplied by \Cref{lem:cylinder-witnesses} and gives
$O(e^{-B^{1+\eps/2}})$.
The additional factor $\mathfrak C_B$ is absorbed by the same estimate, since
\eqref{eq:cylinder-budget} gives $\log \mathfrak C_B=o(B)$ whereas the saving
before this final simplification is $\exp(-\Omega(B^{1+\eps}))$.
This proves the proposition.
\end{proof}

The reduction has preserved a signed average at the unfixed coordinates.
For its use in the next section, it is useful to distinguish
$\mathcal C'=\mathcal C\setminus\mathcal F$ and
$\mathcal Z'=\mathcal Z\setminus\mathcal F$.  For each list,
\begin{equation}\label{eq:cylinder-core-order}
 \left|\E_{\mathcal S\setminus\mathcal F}
              K_{\boldsymbol d}I_{\mathfrak L}\right|
 \le \E_{\mathcal C'}
       \left|\E_{\mathcal Z'}K_{\boldsymbol d}I_{\mathfrak L}\right|.
\end{equation}
This is only Fubini's theorem and the triangle inequality.  In the
inner average all core coordinates are held fixed, so the cutoffs are
fixed there; the free center coordinates still retain their signs.
The core coordinates can subsequently be integrated against positive
bounds.  In particular, \eqref{eq:cylinder-core-order} does not enlarge
$\mathcal F$ or change the private-coordinate rank used in the sieve.
